\documentclass{article}
\usepackage[T1]{fontenc}
\usepackage{lmodern}
\usepackage{graphicx}
\usepackage{amsmath,amssymb}
\usepackage[a4paper,margin=2cm]{geometry}
\usepackage[absolute,overlay]{textpos}
\setlength{\TPHorizModule}{1mm}
\setlength{\TPVertModule}{1mm}
\usepackage[indonesian]{babel}
\usepackage{hyperref}
\setlength{\emergencystretch}{2em}
% unit: Halaman 68

\begin{document}

% === Halaman 68 (llm) ===

\begin{itemize}
    \item \textit{Playfair cipher} termasuk ke dalam kelompok \textit{polygram cipher}
    \item Pada \textit{polygram cipher}, satu blok huruf plainteks disubstitusi dengan satu blok huruf cipherteks.
    \item Jika satu blok panjangnya 2 huruf, maka ia disebut digram (\textit{bigram}), jika 3 huruf disebut trigram, dst. \textit{Playfair cipher} melakukan enkripsi/dekripsi dalam bentuk bigram.

    Misalnya, bigram \texttt{AS} diganti dengan \texttt{RT}, \texttt{BY} diganti dengan \texttt{SL}
\end{itemize}

\end{document}
